% ----------------------
\subsection{Section 87 (25 December 1832)}
% ----------------------
	\label{DC87}
	
	% -----
	\subsubsection{Historical Background}
	% -----
		\label{DC87-HB}
		
		% --
		\paragraph{JSP}
		% --
		
			``On Christmas Day, 25 December 1832, JS dictated this revelation warning of the outbreak of war across all nations, beginning in South Carolina. Remarking on the context of this revelation, a later JS history states: `Appearances of troubles among the nations, became more visible, this season, than they had previously done, since the church began her journey out of the wilderness.' The \textit{Painesville Telegraph} of 21 December 1832 highlighted some of these problems. It contained an article titled `Revenge and Magnanimity. A Tale of the Cholera' about the worldwide cholera epidemic, as well as information about a plague in India that was killing 150 to 200 people a day. The newspaper also included extensive coverage of the passage of a resolution by a Nullification Convention held in November in South Carolina. This resolution declared the federal tariff acts of 1828 and 1832, which levied high duties against imports, `null and void' in the state. Many South Carolina residents believed the acts were passed solely to protect northern manufacturing at the expense of the South. Not only did South Carolinians claim the right to nullify the law, they also stated their willingness to `organize a separate Government' should the federal government try to enforce the tariffs in the state. The governor called for two thousand men to form a militia `for the defence of Charleston and its dependencies.' President Andrew Jackson responded quickly to this resolution, stating, according to the \textit{Telegraph}, `that the laws and the Union must be maintained, at all events.' Because Painesville, Ohio, was only about ten miles from Kirtland, Ohio, it is probable that JS saw or heard about the articles in the 21 December \textit{Telegraph} within a day or so. These developments troubled JS, who saw in them the threat of the `immediate dissolution' of the United States. Indeed, the 25 December revelation predicted that rebellion on the part of South Carolina would lead not only to civil war and war among nations but also to slave rebellions and an uprising of remnants of the house of Israel. This violence, combined with plague and other natural disasters, would ultimately lead to the `full end of all Nations.' Using millenarian language, the revelation cast such events as portents of the return of Jesus Christ to the earth.%
				\footnote{%
					% \label{}%
					HDDC 1104--1109 has a lot more information about the connection to the Civil War, including a bit about what happened after the war broke out, and how people saw this revelation. There's an interesting article from the 5 May 1861 edition of the \textit{Philadelphia Sunday Mercury}, where it quotes this section and comments a bit on it and then finishes with, ``Have we not had a prophet among us?''
					} 
			
			``Frederick G. Williams wrote the revelation as JS dictated it, but the original manuscript is no longer extant. Probably between late January and late February 1833, Williams copied the revelation into Revelation Book 2, titling it `Prophecy given Dec 25--- 1832 concerning concerning the wars' in that book's index.''
			
	% -----
	\subsubsection{Versions}
	% -----
		\label{DC87-V}
		
		See the \href{https://drive.google.com/file/d/1goAvNz8jS4R8slYfMG_db55Ty2z_vmgK/view?usp=sharing}{sections comparison} document.
		
		This revelation was widely circulated in manuscript form for many years before it was first published in 1851. Woodford lists eight manuscript copies made before 1841, and those were just the ones he knew about; there could have been more. It appears as though it was first published in 1851 in the PGP, and then Orson Pratt published it in \textit{The Seer} in 1854. It later appeared a bunch of times in the MS---the one listed below is just the first one. See HDDC 1112--1117.
		
		It first appeared in the DC in 1876.
	
		\begin{compactdesc}
			\item[JSP] \hyperref[EMS-1832-JS-25-DEC-1832]{25 December 1832}
			\item[BC] ---
			\item[DC 1835] ---
			\item[TS] ---
			\item[DC 1844] ---
			\item[PGP 1851] 35
			\item[MS] 22:4 (28 January 1860), 51
		\end{compactdesc}

	% -----
	\subsubsection{Section 87:1} % *DC87-1 
	% -----
		\label{DC87-1}

		VERILY, thus saith the Lord concerning the wars that will shortly come to pass, beginning at the rebellion of South Carolina, which will eventually terminate in the death and misery of many souls;

		% \begin{framed}
		% \scriptsize
			% \begin{compactdesc}
				% \item[JSP] 
				% % \item[BC] 
				% % \item[]
			% \end{compactdesc}
		% \end{framed}

	% -----
	\subsubsection{Section 87:2} % *DC87-2 
	% -----
		\label{DC87-2}

		And the time will come that war will be poured out upon all nations, beginning at this place.

		% \begin{framed}
		% \scriptsize
			% \begin{compactdesc}
				% \item[JSP] 
				% % \item[BC] 
				% % \item[]
			% \end{compactdesc}
		% \end{framed}

	% -----
	\subsubsection{Section 87:3} % *DC87-3 
	% -----
		\label{DC87-3}

		For behold, the Southern States shall be divided against the Northern States, and the Southern States will call on other nations, even the nation of Great Britain,%
			\footnote{%
				% \label{}%
				According to Wikipedia, ``The United Kingdom of Great Britain and Ireland remained officially neutral throughout the American Civil War (1861--1865). It legally recognised the belligerent status of the Confederate States of America (CSA) but never recognised it as a nation and neither signed a treaty with it nor ever exchanged ambassadors.''
				}
		as it is called, and they shall also call upon other nations, in order to defend themselves against other nations; and then war shall be poured out upon all nations.

		% \begin{framed}
		% \scriptsize
			% \begin{compactdesc}
				% \item[JSP] 
				% % \item[BC] 
				% % \item[]
			% \end{compactdesc}
		% \end{framed}

	% -----
	\subsubsection{Section 87:4} % *DC87-4 
	% -----
		\label{DC87-4}

		And it shall come to pass, after many days, slaves shall rise up against their masters, who shall be marshaled and disciplined for war.

		% \begin{framed}
		% \scriptsize
			% \begin{compactdesc}
				% \item[JSP] 
				% % \item[BC] 
				% % \item[]
			% \end{compactdesc}
		% \end{framed}

	% -----
	\subsubsection{Section 87:5} % *DC87-5 
	% -----
		\label{DC87-5}

		And it shall come to pass also that the remnants who are left of the land will marshal themselves, and shall become exceedingly angry, and shall vex the Gentiles with a sore vexation.

		% \begin{framed}
		% \scriptsize
			% \begin{compactdesc}
				% \item[JSP] 
				% % \item[BC] 
				% % \item[]
			% \end{compactdesc}
		% \end{framed}

	% -----
	\subsubsection{Section 87:6} % *DC87-6 
	% -----
		\label{DC87-6}

		And thus, with the sword and by bloodshed the inhabitants of the earth shall mourn; and with famine, and plague, and earthquake, and the thunder of heaven, and the fierce and vivid lightning also, shall the inhabitants of the earth be made to feel the wrath, and indignation, and chastening hand of an Almighty God, until the consumption decreed hath made a full end of all nations;

		% \begin{framed}
		% \scriptsize
			% \begin{compactdesc}
				% \item[JSP] 
				% % \item[BC] 
				% % \item[]
			% \end{compactdesc}
		% \end{framed}

	% -----
	\subsubsection{Section 87:7} % *DC87-7 
	% -----
		\label{DC87-7}

		That the cry of the saints, and of the blood of the saints, shall cease to come up into the ears of the Lord of \hyperref[SABAOTH]{Sabaoth}, from the earth, to be avenged of their enemies.

		% \begin{framed}
		% \scriptsize
			% \begin{compactdesc}
				% \item[JSP] 
				% % \item[BC] 
				% % \item[]
			% \end{compactdesc}
		% \end{framed}

	% -----
	\subsubsection{Section 87:8} % *DC87-8 
	% -----
		\label{DC87-8}

		Wherefore, stand ye in holy places, and be not moved, until the day of the Lord come; for behold, it cometh quickly, saith the Lord.  Amen.

		% \begin{framed}
		% \scriptsize
			% \begin{compactdesc}
				% \item[JSP] 
				% % \item[BC] 
				% % \item[]
			% \end{compactdesc}
		% \end{framed}